\documentclass[12pt]{article}
\usepackage[utf8]{inputenc}
\usepackage[T1]{fontenc}
\usepackage{hyperref}
\title{Zhuang, Eastern Hongshuihe Language Resources for K-12 Curriculum Development}
\author{Compiled by Claude Sonnet 5.5 (claude-sonnet-5-5) for Andr\'as Kornai}
\date{}
\begin{document}
\maketitle

\section{Introduction}
Eastern Hongshuihe Zhuang (ISO 639-3 zeh) is a Northern Tai language of the Kra-Dai family and belongs to the Zhuang macrolanguage. It is spoken in western Guangxi and Wenshan Prefecture, Yunnan, China. Wikipedia gives about 1.2 million speakers (2007). No materials specific to this variety were found, so the resources below concern Zhuang generally.

\section{Primary Educational Resources}
No textbook titles or publishers were identified in the searches. The sources report that Standard Zhuang is taught in some bilingual primary schools in Guangxi (108 such schools in a 2013 report), but only a small share of schools teach written Zhuang. Background on the standard script and the bilingual policy is at \url{https://en.wikipedia.org/wiki/Standard_Zhuang}. A 2013 blog post discusses the Zhuang proficiency exam (VSSG), which was promoted to support bilingual Zhuang-Chinese education: \url{https://bruce-humes.com/2013/12/15/candidates-for-romanized-zhuang-exam-few-and-far-between/}

\section{Linguistic and Reference Materials}
\begin{itemize}
\item Eastern Hongshuihe Zhuang, Wikipedia: \url{https://en.wikipedia.org/wiki/Eastern_Hongshuihe_Zhuang}
\item Zhuang languages, Wikipedia: \url{https://en.wikipedia.org/wiki/Zhuang_languages}
\item Zhuang writing systems, Omniglot (Latin-based Standard Zhuang and the traditional Sawndip characters): \url{https://omniglot.com/writing/zhuang.htm}
\end{itemize}

\section{Religious and Community Resources}
Joshua Project lists Bible translation for Eastern Hongshuihe Zhuang as started: \url{https://joshuaproject.net/people_groups/19803}

\section{Digital and Audio Resources}
Global Recordings Network has a language page for Eastern Hongshuihe Zhuang: \url{https://globalrecordings.net/en/language/zeh}. The Wikitongues language page is at \url{https://wikitongues.org/languages/zeh/}. Neither page was opened, so their contents are not confirmed here.

\section{Curriculum Development Context}
The language is vigorous (EGIDS 6a per Wikitongues). Standard Zhuang, based on the Wuming dialect, exists as a written standard and is used in some bilingual schools, so a curriculum could adapt it. Variety-specific adaptation of Eastern Hongshuihe material would still be needed. Locating actual Zhuang textbooks, for example through Guangxi educational publishers, is a next step that the searches did not complete.
\end{document}